\section{行列の指数関数と時間発展}
\label{sec:matrix-exponential}

量子位相推定の入力はユニタリ写像である．
一方，HHLで扱いたい情報はエルミート行列$A$の固有値である．
そこで本節では，行列の指数関数を使って$A$からユニタリ写像を作る
\cite{shimada2020}．

\begin{defi}[行列指数関数]
正方行列$B\in\mathbb C^{N\times N}$に対し，
\begin{equation}
 e^B:=\sum_{k=0}^{\infty}\frac{B^k}{k!}
 \label{eq:matrix-exponential-definition}
\end{equation}
を$B$の指数関数という．
\end{defi}

\begin{prop}
式\eqref{eq:matrix-exponential-definition}の級数は任意の正方行列$B$について収束する．
\end{prop}
\begin{proof}
積に対して$\lVert CD\rVert\leq\lVert C\rVert\lVert D\rVert$を満たす行列ノルムを用いる．
すると$\lVert B^k\rVert\leq\lVert B\rVert^k$であるから，
\[
 \sum_{k=0}^{\infty}\left\lVert\frac{B^k}{k!}\right\rVert
 \leq\sum_{k=0}^{\infty}\frac{\lVert B\rVert^k}{k!}
 =e^{\lVert B\rVert}<\infty
\]
右辺は実数の指数関数の級数であり，有限の値に収束する．
したがって，行列の級数も絶対収束する．
\end{proof}

\begin{theo}
行列$A$をエルミートとし，スペクトル分解を
\begin{equation}
 A=\sum_j\lambda_j\ket{u_j}\bra{u_j}
 \label{eq:hermitian-spectral-decomposition}
\end{equation}
とする．このとき任意の実数$t$について
\begin{equation}
 e^{\mathrm{i}At}
 =\sum_j e^{\mathrm{i}\lambda_jt}\ket{u_j}\bra{u_j}
 \label{eq:matrix-exponential-spectral}
\end{equation}
であり，$e^{\mathrm{i}At}$はユニタリである．
\end{theo}
\begin{proof}
正規直交性より
$(\ket{u_j}\bra{u_j})(\ket{u_l}\bra{u_l})
=\delta_{jl}\ket{u_j}\bra{u_j}$である．したがって帰納法により
\[
 A^k=\sum_j\lambda_j^k\ket{u_j}\bra{u_j}
\]
を得る．これを式\eqref{eq:matrix-exponential-definition}へ代入し，有限個の$j$と
絶対収束する級数の和を交換すると
\[
 e^{\mathrm{i}At}
 =\sum_j\left(\sum_{k=0}^{\infty}
   \frac{(\mathrm{i}\lambda_jt)^k}{k!}\right)
   \ket{u_j}\bra{u_j}
 =\sum_je^{\mathrm{i}\lambda_jt}\ket{u_j}\bra{u_j}
\]
となる．また$\lambda_j,t$は実数なので
\begin{align*}
 (e^{\mathrm{i}At})^\dagger e^{\mathrm{i}At}
 &=\sum_{j,l}e^{-\mathrm{i}\lambda_jt}e^{\mathrm{i}\lambda_lt}
   \ket{u_j}\braket{u_j|u_l}\bra{u_l}\\
 &=\sum_j\ket{u_j}\bra{u_j}=I
\end{align*}
である．逆順の積も同様であるからユニタリである．
\end{proof}

式\eqref{eq:matrix-exponential-spectral}から，各固有ベクトルについて
\begin{equation}
 e^{\mathrm{i}At}\ket{u_j}
 =e^{\mathrm{i}\lambda_jt}\ket{u_j}
 =e^{2\pi\mathrm{i}\varphi_j}\ket{u_j},
 \qquad
 \varphi_j=\frac{\lambda_jt}{2\pi}\pmod 1
 \label{eq:matrix-eigenvalue-to-phase}
\end{equation}
が成り立つ．したがって$e^{\mathrm{i}At}$に量子位相推定を適用すれば，
位相$\varphi_j$を通して固有値$\lambda_j$を推定できる．

ただし，位相は1を法としてしか区別できない．固有値が既知の区間
$[\lambda_{\min},\lambda_{\max}]$にある場合，異なる固有値が同じ位相へ
折り返されないよう$t$を選び，位相から固有値へ戻す規約を定める必要がある．
負の固有値を含む場合には，位相レジスタ上の符号表現も必要となる．

本論文では，制御された$e^{\mathrm{i}At}$と，位相推定で使うそのべき乗を
必要な精度で実行できると仮定する．
行列$A$からこの量子回路を作る処理をハミルトニアンシミュレーションという．
本論文ではその具体的な回路までは扱わない．
ただし，必要な計算量は第4章でHHLの条件として書く．
